\documentclass{article}
%  ╔══════════════════════════════════════════════════════════╗
%  ║                         Title                            ║
%  ╚══════════════════════════════════════════════════════════╝
%NOTE: I want \title,\author to be near top, and this \renewcommand{\maketitle} fails if these are above it.
\usepackage[utf8]{inputenc}         % Used to compile any UTF-8 Character (and supports Unicode)
\usepackage{titling}

\renewcommand{\maketitle}{
    \quad
    \vspace{1em}
  \begin{center}
      \LARGE\bfseries \thetitle \par
      \vspace{0.6em}
      \large
  \end{center}
    \vspace{1em}
}

\title{Research in Algebraic Geometry}
\author{Nathanael Chwojko-Srawley}
\date{\today}

%  ╔══════════════════════════════════════════════════════════╗
%  ║                         packages                         ║
%  ╚══════════════════════════════════════════════════════════╝

\usepackage{extarrows}         % Used to compile any UTF-8 Character (and supports Unicode)
\usepackage{stackengine}
\usepackage{colortbl}
\usepackage{scalerel}
\usepackage{amsmath, amssymb} % Used to write better mathematical expressions (ex. \mathbb{R})
\usepackage{stackrel}         % for left-right arrows, staking symbol above and bellow
\usepackage{mathrsfs}
\usepackage{bbm}                     % for mathbbm{1}
\usepackage{euscript}
\usepackage{commutative-diagrams}
\usepackage[font=itshape]{quoting}
\usepackage{mathtools}
  \usepackage{enumitem}
\usepackage{amsthm}           % Used to create theorem enviroments.
\usepackage{graphicx}         % Let's you use images in LaTeX; ex the \includegraphics command.
\usepackage{hhline}
%\usepackage{luatexja}
%\usepackage{fontspec}
\usepackage{dsfont}
\usepackage{wasysym}          % emoji's!
\usepackage{float}            % package with ``H'' for figures and tables, ex. Images, tables, etc
%\usepackage{subcaption}      % More options with sub-captions.
\usepackage{setspace}         % More flexibility on spacing in document.
\usepackage{hyperref}         % let's me put hyperlinks
%\usepackage{tikz}             % for commutative diagrams https://tex.stackexchange.com/questions/419221/how-to-do-the-pushout-with-universal-property
\usepackage{tikz-cd}          % for commutative diagrams https://tex.stackexchange.com/questions/419221/how-to-do-the-pushout-with-universal-property
%\usepackage{quiver}           % for better commutative diagrams
\usepackage{bookmark}         % creating heading shortcut bar on the left as in word
\usepackage{longtable}        % create tables that extend past one page
%\usepackage{siunitx}         % required for alignment
%\usepackage{multirow}        % for multiple rows
\usepackage{microtype}        % makes nice text. Fixes some under-/over- filled box problems.
%\usepackage{fancyvrb}        % Let's you write code (different style than listings).
\usepackage[most]{tcolorbox} 	      % Makes nice boxes for thms, cor, prop, or anything else.
%\usepackage{enumitem}        % More flexibility on enumeration (don't know much).
\usepackage{xcolor}           % Colour text.
\usepackage{wrapfig}          % To wrap text around figure
\usepackage{listings}         % Create nice colour code blocks (customized in preamble)
\usepackage{thmtools}         % Customize theorem environment (in preamble).
%\usepackage{marvosym}        % Provide ``basic'' symbols (hazard, religious, smiley face etc.)
\usepackage{array}            % \begin{table}{>{\em}c c} -> 1st column italic (bfseries for bold)
%\usepackage[english]{babel}  % If I want to blindtext in english.
\usepackage{blindtext}
\usepackage[a4paper, total={6in, 8.5in}]{geometry}
%\usepackage[symbol]{footmisc} % For daggers and asterisk for footnotes. 1: *, 2: dagger, etc.
\usepackage{fancyhdr}
\usepackage{titlesec}
%\usepackage{pst-plot}          %To make nice plots: https://tex.stackexchange.com/questions/47594/how-can-i-make-a-graph-of-a-function#47596
% \usepackage[answerdelayed]{exercise}
\usepackage{etoolbox}          % for Dynkin diagrams
\usepackage{dynkin-diagrams}   % for Dynkin diagrams

% ╔═════════════════════════════════════════════════════╗
% ║           for figures via inkspace                  ║
% ╚═════════════════════════════════════════════════════╝

% to import figures via: https://github.com/gillescastel/inkscape-figures
\usepackage{import}
\usepackage{pdfpages}
\usepackage{transparent}


% This command is the ONLY command imported via the packages snippet, think of it as a tiny package written out
\newcommand{\incfig}[2][1]{%
    \def\svgwidth{#1\columnwidth}
    \import{./figures/}{#2.pdf_tex}
}

% ╔═══════════════════════════════════════════════════════════╗
% ║                 bibliography and citing                   ║
% ╚═══════════════════════════════════════════════════════════╝

\usepackage{csquotes}
% \usepackage[backend=biber,citestyle=alphabetic,bibstyle=authortitle]{biblatex}
\usepackage{biblatex}

\addbibresource{bibliography.bib}

% ╔═══════════════════════════════════════════════════════════╗
% ║                    colours (colors)                       ║
% ╚═══════════════════════════════════════════════════════════╝

\definecolor{dkgreen}{rgb}{0,0.6,0}
\definecolor{gray}{rgb}{0.5,0.5,0.5}
\definecolor{mauve}{rgb}{0.58,0,0.82}
\definecolor{lightyellow}{rgb}{1,1,0.6}
\definecolor{reddish}{HTML}{A01A3A}

%  ╔══════════════════════════════════════════════════════════╗
%  ║                       book format                        ║
%  ╚══════════════════════════════════════════════════════════╝

\pagestyle{fancy}

\setlength\headheight{20pt}

\renewcommand{\sectionmark}[1]{\markright{\thesection.\ #1}}

\fancyfoot[C,CO]{\textcolor{gray}{Nathanael Chwojko-Srawley}}
\fancyfoot[R,RO]{\thepage}{}

% Create a new page style for the first page
\fancypagestyle{firstpage}{
  \fancyhf{} % Clear header and footer
  \fancyhead[L]{\theauthor}
  \fancyhead[R]{\today} % Date in the top-right
  \fancyfoot[C]{\textcolor{gray}{Nathanael Chwojko-Srawley}}
  \fancyfoot[R]{\thepage}
}

% Apply the first page style conditionally
\AtBeginDocument{\thispagestyle{firstpage}}

%have format the title command
\newcommand{\chapter}[1]{%
  \clearpage
  \vspace*{30pt} % Space before the chapter title
  {\normalfont\huge\bfseries #1\par} % Chapter title formatting
  \vspace{20pt} % Space after the chapter title
}


\setlength{\parindent}{0pt} % don't like indentation infront of paragraphs
\setlength{\parskip}{1ex}   %so that there is still space between paragarphs


%  ╔══════════════════════════════════════════════════════════╗
%  ║                       Shortcuts                          ║
%  ╚══════════════════════════════════════════════════════════╝

%\ExerciseLevelInToc

% I like original mathcal better, but need lowercase.
\DeclareFontFamily{OT1}{pzc}{}
\DeclareFontShape{OT1}{pzc}{m}{it}{<-> s * [1.10] pzcmi7t}{}
\DeclareMathAlphabet{\mathpzc}{OT1}{pzc}{m}{it}


% mathbb
\newcommand{\A}{{\mathbb{A}}}
\newcommand{\C}{{\mathbb{C}}}
\newcommand{\F}{{\mathbb{F}}}
\newcommand{\N}{{\mathbb{N}}}
\newcommand{\Q}{{\mathbb{Q}}}
\newcommand{\R}{{\mathbb{R}}}
\newcommand{\V}{{\mathbb{V}}}
\newcommand{\Z}{{\mathbb{Z}}}
\newcommand{\BI}{{\mathbb{I}}}
\renewcommand{\H}{{\mathbb{H}}}
\renewcommand{\P}{{\mathbb{P}}}


% mathcal
% \newcommand{\co}{{\mathfrak{o}}} %mathcal{o} looks like wreath product, so I changed it to mathfrak
\newcommand{\co}{{\mathpzc{o}}}
\newcommand{\CA}{{\mathcal{A}}}
\newcommand{\CB}{{\mathcal{B}}}
\newcommand{\CC}{{\mathcal{C}}}
\newcommand{\CD}{{\mathcal{D}}}
\newcommand{\CF}{{\mathcal{F}}}
\newcommand{\CG}{{\mathcal{G}}}
\newcommand{\CH}{{\mathcal{H}}}
\newcommand{\CI}{{\mathcal{I}}}
\newcommand{\CK}{{\mathcal{K}}}
\newcommand{\CL}{{\mathcal{L}}}
\newcommand{\CM}{{\mathcal{M}}}
\newcommand{\CN}{{\mathcal{N}}}
\newcommand{\CO}{{\mathcal{O}}}
\newcommand{\CQ}{{\mathcal{Q}}}
\newcommand{\CR}{{\mathcal{R}}}
\newcommand{\CS}{{\mathcal{S}}}
\newcommand{\CT}{{\mathcal{T}}}
\newcommand{\CV}{{\mathcal{V}}}
\newcommand{\CZ}{{\mathcal{Z}}}
% EXCPETION:
\newcommand{\CP}{{\mathcal{P}}}
% \newcommand{\CP}{\mathbb{C}\mathrm{P}}
\newcommand{\RP}{\mathbb{R}\mathrm{P}}


%Euscript
\newcommand{\EB}{\EuScript{B}}
\newcommand{\EC}{{\EuScript{C}}}
\newcommand{\EF}{{\EuScript{F}}}
\newcommand{\EL}{{\EuScript{L}}}
\newcommand{\EO}{{\EuScript{O}}}


%mathfrak (lowercase)
\newcommand{\fa}{{\mathfrak{a}}}
\newcommand{\fb}{{\mathfrak{b}}}
\newcommand{\fc}{{\mathfrak{c}}}
\newcommand{\fd}{{\mathfrak{d}}}
\newcommand{\fm}{{\mathfrak{m}}}
\newcommand{\fn}{{\mathfrak{n}}}
\newcommand{\fo}{{\mathfrak{o}}}
\newcommand{\fp}{{\mathfrak{p}}}
\newcommand{\fq}{{\mathfrak{q}}}


%mathfrak (upper case)
\newcommand{\FA}{{\mathfrak{A}}}
\newcommand{\FB}{{\mathfrak{B}}}
\newcommand{\FC}{{\mathfrak{C}}}
\newcommand{\FD}{{\mathfrak{D}}}
\newcommand{\FK}{{\mathfrak{K}}}
\newcommand{\FM}{{\mathfrak{M}}}
\newcommand{\FN}{{\mathfrak{N}}}
\newcommand{\FO}{{\mathfrak{O}}}
\newcommand{\FP}{{\mathfrak{P}}}


%mathscr
\newcommand{\ff}{{\mathscr{F}}}
\newcommand{\SA}{\mathscr{A}}
\newcommand{\SC}{\mathscr{C}}
\newcommand{\SF}{\mathscr{F}}
\newcommand{\SG}{\mathscr{G}}
\newcommand{\SH}{\mathscr{H}}
\newcommand{\SI}{\mathscr{I}}
\newcommand{\SK}{\mathscr{K}}
\newcommand{\SN}{\mathscr{N}}
\newcommand{\SQ}{\mathscr{Q}}
\newcommand{\SR}{\mathscr{R}}
\newcommand{\SV}{\mathscr{V}}
\newcommand{\SZ}{\mathscr{Z}}


% operatorname -- 2 letters (lowercase and upper case)
\newcommand{\ab}{{\operatorname{ab}}}
\newcommand{\ad}{{\operatorname{ad}}}
\newcommand{\ch}{{\operatorname{ch}}}
\newcommand{\ev}{{\operatorname{ev}}}
\newcommand{\id}{{\operatorname{id}}}
\newcommand{\im}{{\operatorname{im}}}
\newcommand{\nm}{{\operatorname{Nm}}}
\newcommand{\op}{{\operatorname{op}}}
\newcommand{\rk}{{\operatorname{rk}}}
\newcommand{\sh}{{\operatorname{sh}}}
\newcommand{\tr}{{\operatorname{tr}}}

\newcommand{\Ab}{{\operatorname{Ab}}}
\newcommand{\Br}{{\operatorname{Br}}}
\newcommand{\Ch}{{\operatorname{Ch}}}
\newcommand{\Cl}{{\operatorname{Cl}}}
\newcommand{\GL}{{\operatorname{GL}}}
\newcommand{\Nm}{{\operatorname{Nm}}}
\newcommand{\SL}{{\operatorname{SL}}}

\renewcommand{\Re}{{\operatorname{Re}}}
\renewcommand{\Im}{{\operatorname{Im}}}


%categories
\newcommand{\Grp}{{\textbf{Grp}}}
\newcommand{\Fld}{{\textbf{Fld}}}
\newcommand{\Sh}{{\textbf{Sh}}}
\newcommand{\Set}{{\textbf{Set}}}
\newcommand{\Mod}{{\textbf{Mod}}}
\newcommand{\PSh}{{\textbf{PSh}}}
\newcommand{\Sch}{{\textbf{Sch}}}
\newcommand{\Top}{{\textbf{Top}}}
\newcommand{\RgSp}{{\textbf{RgSp}}}
\newcommand{\Ring}{{\textbf{Ring}}}


%operatorname -- 3+ letters (lowercase and upper case)
\newcommand{\rad}{{\operatorname{rad}}}
\newcommand{\sep}{{\operatorname{sep}}}
\newcommand{\res}{{\operatorname{res}}}
\newcommand{\sgn}{{\operatorname{sgn}}}
\newcommand{\pre}{{\operatorname{pre}}}
\newcommand{\ord}{{\operatorname{ord}}}
\newcommand{\vol}{{\operatorname{vol}}}
\newcommand{\lcm}{{\operatorname{lcm}}}

\newcommand{\conj}{{\operatorname{conj}}}
\newcommand{\disc}{{\operatorname{disc}}}
\newcommand{\stab}{{\operatorname{stab}}}
\newcommand{\kdim}{{\operatorname{kdim}}}
\newcommand{\diag}{{\operatorname{diag}}}

\newcommand{\trdeg}{\operatorname{trdeg}}
\newcommand{\coker}{{\operatorname{coker}}}
\newcommand{\codim}{{\operatorname{codim}}}

\newcommand{\Ann}{{\operatorname{Ann}}}
\newcommand{\Aut}{{\operatorname{Aut}}}
\newcommand{\Der}{{\operatorname{Der}}}
\newcommand{\Div}{{\operatorname{Div}}}
\newcommand{\Emb}{{\operatorname{Emb}}}
\newcommand{\End}{{\operatorname{End}}}
\newcommand{\Ext}{{\operatorname{Ext}}}
\newcommand{\Gal}{{\operatorname{Gal}}}
\newcommand{\Hom}{{\operatorname{Hom}}}
\newcommand{\Ind}{{\operatorname{Ind}}}
\newcommand{\Inn}{{\operatorname{Inn}}}
\newcommand{\Int}{{\operatorname{int}}}
\newcommand{\Irr}{{\operatorname{Irr}}}
\newcommand{\Jac}{{\operatorname{Jac}}}
\newcommand{\Mor}[1]{\operatorname{Mor}(\textbf{#1})}
\newcommand{\Nil}{{\operatorname{Nil}}}
\newcommand{\Obj}{{\operatorname{Obj}}}
\newcommand{\Out}{{\operatorname{Out}}}
\newcommand{\Pic}{{\operatorname{Pic}\,}}
\newcommand{\Rep}{{\operatorname{Rep}}}
\newcommand{\Res}{{\operatorname{Res}}}
\newcommand{\Spl}{{\operatorname{Spl}}}
\newcommand{\Syl}{{\operatorname{Syl}}}
\newcommand{\Sym}{{\operatorname{Sym}}}
\newcommand{\Tor}{{\operatorname{Tor}}}

\newcommand{\Char}{{\operatorname{char}}}
\newcommand{\Frac}{{\operatorname{Frac}}}
\newcommand{\Frob}{{\operatorname{Frob}}}
\newcommand{\Proj}{{\operatorname{Proj}\,}}
\newcommand{\RMod}{{}_R\operatorname{Mod}}
\newcommand{\SExt}{{\mathcal{E}\emph{xt}}}
\newcommand{\SHom}{{\emph{Hom}}}
\newcommand{\Span}{{\operatorname{span}}}
\newcommand{\Spec}{{\operatorname{Spec}\,}}
\newcommand{\Supp}{{\operatorname{Supp}}}
\newcommand{\Tens}[1]{#1\otimes --}
\newcommand{\fHom}[1]{\operatorname{Hom}(#1, --)}

\newcommand{\mSpec}{{\operatorname{mSpec}}}
\newcommand{\cofHom}[1]{\operatorname{Hom}(--, #1)}


% connectors
% \renewcommand{\mid}{\big|}
\newcommand{\sse}{\subseteq}
\newcommand{\subg}{\leqslant}
\newcommand{\supg}{\geqslant}
\newcommand{\ssne}{\subsetneq}
\newcommand{\isosubg}{\lesssim}
\newcommand{\nsse}{\not\subseteq}
\newcommand{\nsubg}{\trianglelefteq}
\newcommand{\nsupg}{\trianglerighteq}
\newcommand{\csubg}{\blacktriangleleft}
\renewcommand{\mid}{\big|}


% Arrows
\newcommand{\rw}{\rightarrow}
\newcommand{\Rw}{\Rightarrow}
\newcommand{\lw}{\leftarrow}
\newcommand{\Lw}{\Leftarrow}
\newcommand{\lrw}{\leftrightarrow}
\newcommand{\LRw}{\Leftrightarrow}
\newcommand{\acton}{\curvearrowright}
\newcommand{\actson}{\curvearrowright}
\newcommand\mapsfrom{\mathrel{\reflectbox{\ensuremath{\mapsto}}}}

% arrows for commutative diagram: create four new angled double-struck arrows
\newcommand\myrot[1]{\mathrel{\rotatebox[origin=c]{#1}{$\Rightarrow$}}}
\newcommand\NEarrow{\myrot{45}}
\newcommand\NWarrow{\myrot{135}}
\newcommand\SWarrow{\myrot{-135}}
\newcommand\SEarrow{\myrot{-45}}


%shorthands
\newcommand{\oln}[1]{{\overline{#1}}}
\renewcommand{\bf}[1]{\textbf{#1}}
\newcommand{\qn}{\quad\newline}


% limits
\newcommand{\Lim}{{\varprojlim}}
\newcommand{\coLim}{{\varinjlim}}
\newcommand{\Dlim}{\lim\limits_{\longrightarrow}}
\newcommand{\coDlim}{\lim\limits_{\longleftarrow}}


%for saturated ideals in the localization section (I don't like these, I think I'll delete)
\newcommand{\LIm}{{}^e }
\newcommand{\LPre}{{}^c }

%  ╔══════════════════════════════════════════════════════════╗
%  ║                    Custom Commands                       ║
%  ╚══════════════════════════════════════════════════════════╝

% swap varphi and phi
\let\temp\phi
\let\phi\varphi
\let\varphi\temp

\newcommand{\placeholder}{\_\_\_\_\_}

\newcommand{\set}[2]{\left\{#1 \ \middle|\ #2\right\}}

%mod should not have the crazy amount of space to its right!
\renewcommand{\mod}{\,\operatorname{mod}\,}

% dcups
\newcommand{\dcup}{\sqcup}
\newcommand{\Dcup}{\coprod}
% \newcommand{\Dcup}{\bigsqcup}

\newcommand{\nbot}{\mathrel{\rlap{$\bot$}\mkern2mu/}}

%a nice large circle
\newcommand{\tikzcircle}[2][red,fill=black]{\tikz[baseline=-0.5ex]\draw[#1,radius=#2] (0,0) circle ;}%

\makeatletter
\newcommand*\bigcdot{\mathpalette\bigcdot@{.5}}
\newcommand*\bigcdot@[2]{\mathbin{\vcenter{\hbox{\scalebox{#2}{$\m@th#1\bullet$}}}}}
\makeatother

%somehow not a default command
\def\rddots{\mathstrut^{.^{.^{.}}}}

%% new delimiter
\DeclarePairedDelimiter{\ceil}{\lceil}{\rceil}


%  ╔══════════════════════════════════════════════════════════╗
%  ║                     environments                         ║
%  ╚══════════════════════════════════════════════════════════╝


%remember the option: breakable

% create tcolor boxes
\newtcbtheorem[number within=section]{thm}{Theorem}%
{
colback=green!5,
colframe=green!35!black,
fonttitle=\bfseries,
label type=theorem
}{th}

\newtcbtheorem[number within=section, use counter from=thm]{lem}{Lemma}%
{
colback=blue!5,
colframe=black!35!black,
fonttitle=\bfseries,
label type=lemma
}{lm}
\newtcbtheorem[number within=section, use counter from=thm]{prop}{Proposition}%
{
colback=white!5,
colframe=white!35!black,
fonttitle=\bfseries,
label type=proposition
}{pr}
\newtcbtheorem[number within=section, use counter from=thm]{cor}{Corollary}%
{colback=blue!5,
colframe=blue!35!black,
fonttitle=\bfseries,
label type=corollary
}{co}
\newtcbtheorem[number within=section, use counter from=thm]{axiom}{Axiom}%
{colback=black!5,
colframe=yellow!35!black,
fonttitle=\bfseries,
label type=axiom
}{ax}
\newtcbtheorem[number within=section, use counter from=thm]{defn}{Definition}%
{colback=red!5,
colframe=red!35!black,
fonttitle=\bfseries,
label type=definition
}{df}


% for <s-cr> to create \item[$\LRw$] instead of \item
\newenvironment{equivEnumerate}
 {\begin{enumerate}
	}{
\end{enumerate}}



%Custom enviroments
 \newenvironment{note}
 {\begin{tcolorbox}[enhanced, sharp corners, colback=white , borderline={0.2pt}{0pt}{black}]
   \begin{quote}\textbf{Note}:
   }{
   \end{quote}\end{tcolorbox}}

 \newenvironment{intuition}
 {\begin{quote}\textbf{Intuition}:
   }{
   \end{quote}}

\newtcbtheorem[number within=section, use counter from=thm]{titledBox}{}{%
  enhanced,
  sharp corners,
  colback=white,
  colframe=black,
  borderline={0.2pt}{0pt}{black},
  left=2mm,
  right=2mm,
  top=2mm,
  bottom=2mm,
  title style={opacity=1},
  attach title to upper,
  before upper={\textbf{\tcbtitletext}\quad},
  coltitle=black,
  fonttitle=\bfseries,
  separator sign={\quad}
}{box}


%better example and proof environments
\def\exampletext{Example} % If English
\colorlet{colexam}{red!55!black} % Global example color


\newtcbtheorem[number within=section, use counter from=thm]{example}{Example}{% \newtcolorbox[use counter=testexample]{testexamplebox}{%
% Example Frame Start
empty,% Empty previously set parameters
title={\exampletext \thetcbcounter #1},% use \thetcbcounter to access the testexample counter text
% Attaching a box requires an overlay
attach boxed title to top left,
   % Ensures proper line breaking in longer titles
   minipage boxed title,
% (boxed title style requires an overlay)
boxed title style={empty,size=minimal,toprule=0pt,top=4pt,left=3mm,overlay={}},
coltitle=colexam,fonttitle=\bfseries,
before=\par\medskip\noindent,parbox=false,boxsep=0pt,left=3mm,right=0mm,top=2pt,breakable,pad at break=0mm,
   before upper=\csname @totalleftmargin\endcsname0pt, % Use instead of parbox=true. This ensures parskip is inherited by box.
% Handles box when it exists on one page only
overlay unbroken={\draw[colexam,line width=.5pt] ([xshift=-0pt]title.north west) -- ([xshift=-0pt]frame.south west); },
% Handles multipage box: first page
overlay first={\draw[colexam,line width=.5pt] ([xshift=-0pt]title.north west) -- ([xshift=-0pt]frame.south west); },
% Handles multipage box: middle page
overlay middle={\draw[colexam,line width=.5pt] ([xshift=-0pt]frame.north west) -- ([xshift=-0pt]frame.south west); },
% Handles multipage box: last page
overlay last={\draw[colexam,line width=.5pt] ([xshift=-0pt]frame.north west) -- ([xshift=-0pt]frame.south west); }%
}{ex}



\def\prooftext{Proof} % If English
\NewDocumentEnvironment{Proof}{ O{} }
{
    \colorlet{colexam}{black!55!black} % Global example color
    \newtcolorbox[number within=section]{testproofbox}{%
        empty,% Empty previously set parameters
        title={\emph{\prooftext} \thetcbcounter: #1},
        attach boxed title to top left,
        minipage boxed title,
        boxed title style={empty,size=minimal,toprule=0pt,top=4pt,left=3mm,overlay={}},
        coltitle=colexam,
        fonttitle=\bfseries,
        before=\par\medskip\noindent,
        parbox=false,
        boxsep=0pt,
        left=3mm,
        right=0mm,
        top=2pt,
        breakable,
        pad at break=0mm,
        before upper=\csname @totalleftmargin\endcsname0pt,
        % Removed height constraints
        % Added flexible width
        width=\dimexpr\textwidth-3mm\relax,
        % Modified overlays
        overlay unbroken={\draw[colexam,line width=.5pt] ([xshift=-0pt]title.north west) -- ([xshift=-0pt]frame.south west); },
        overlay first={\draw[colexam,line width=.5pt] ([xshift=-0pt]title.north west) -- ([xshift=-0pt]frame.south west); },
        overlay middle={\draw[colexam,line width=.5pt] ([xshift=-0pt]frame.north west) -- ([xshift=-0pt]frame.south west); },
        overlay last={\draw[colexam,line width=.5pt] ([xshift=-0pt]frame.north west) -- ([xshift=-0pt]frame.south west); },%
    }
    \begin{testproofbox}
}
{\end{testproofbox}\endlist}


%  ╔══════════════════════════════════════════════════════════╗
%  ║                    for Dynkin diagrams                   ║
%  ╚══════════════════════════════════════════════════════════╝

\def\row#1/#2!{#1_{\IfStrEq{#2}{}{n}{#2}} & \dynkin{#1}{#2}\\}
\newcommand{\tble}[1]{
   \renewcommand*\do[1]{\row##1!}
   \[
      \begin{array}{ll}\docsvlist{#1}\end{array}
   \]
}

%  ╔══════════════════════════════════════════════════════════╗
%  ║                         links                            ║
%  ╚══════════════════════════════════════════════════════════╝

\definecolor{LinkColour}{HTML}{A64A3B} % favourite
% \definecolor{LinkColour}{HTML}{8C3D32} % 2nd place
\hypersetup{
  colorlinks,
  citecolor=black,
  filecolor=magenta,
  linkcolor=LinkColour,
  urlcolor=blue,
}


%  ╔══════════════════════════════════════════════════════════╗
%  ║                          misc                            ║
%  ╚══════════════════════════════════════════════════════════╝

\stackMath %to be able to stack text in math mode
\makeindex

%import options: all, bibliography, book-formatting, booksSetting, customEnvironments, dynkin, hw-formatting, language-setup, newcommands, packages, preamble, proofs, settings, tcolorbox, test.svg,
\begin{document}
\maketitle

This is a snippet from the end of \emph{EYNTKA Algebra}~\cite{NateAlgebra} that covers areas of research in algebra that a reader
may want to explore. Thus, this blog assumes comfort with anything from that book.


It must be pointed out that the topics covered here are \emph{not} a comprehensive list of areas of research,
simply large ones. Other areas that are not covered in detail here that have active research include, but are
not limited to, the model theory of algebraically and differentially closed fields, probabilistic linear
optimization, mathematics of AI (which can be thought of as a continuation of linear algebra and functional
analysis), equivalent theory, P vs. NP as an Invariance Theory problem, and so forth.

\subsection*{Algebraic Geometry}

The starting point of modern algebraic geometry is
the \emph{affine scheme} which generalizes varieties
by defining the geometric counter-part to number rings, rings with nilpotents or zero-divisors, rings that are
not finitely generated, and overall any commutative ring. These shall be the ``model spaces'' that are glued
together to create \emph{schemes}, where ``glued together'' is geometrically understood as forming equivalence
classes between the points that are being glued and is algebraically understood as adding compatibility
relations between functions (essentially, think of the example of projective space where $f \in \Hom(U_1,
R),g\in \Hom(U_2, R)$ such that $f(z) = g(1/z)$ on $U_1\cap U_2$). A lot of the build-up in modern algebra
geometry is translating the concepts of differential geometry over $\R$ or $\C$ (ex. smoothness, dimensions,
irreducibility, compactness, Hausdorffness, Poincar\'e duality, monodromy, Riemann-Roch) into the non-analytic
setting of schemes. This translation abstracts many of the notions that will require a good deal of build-up,
but the reward is a unified language to deal with number theory and varieties which have many of the geometric
intuitions one would expect. A lot of similarities shall be shared between differential geometry and scheme
theory, but important differences shall also emerge (ex. points containing other points, existence of distinct
points that are arbitrarily closed to each other, functions that are not determined by their value at every
point).

Scheme are the starting point for the study of many different modern fields, but it can be studied in
it of itself. For example to find classification of different types of schemes; this research often involves the use of
\emph{moduli spaces} and the theory of \emph{stacks}.

Another very important problem in Algebraic Geometry is the \emph{Hodge Conjecture} which roughly says says
that homological information about compact complex manifolds $X$ (which from an algebraic view point, are objects
that very closed resemble smooth projective varieties and often are categorically equivalent due to
GAGA~\cite{hallGAGATheorems2022}) can be given by nice subvarieties of $X$. In particular, without defining
any of our terms, it asks which cohomology classes of $H^{k,k}(X)$ can be represented by closed subvarieties
$Z$ where $X$ is given the natural K\"ahler manifold structure. The interested reader wanting to continue the
EYTNKA series can read \emph{K-theory} for an introduction to characteristic classes~\cite{NateKTheory},
\emph{Algebraic Topology} for the cohomology background~\cite{NateAlgTop}, \emph{Algebraic Geometry} for the
algebraic background~\cite{NateAlgGeo}, and \emph{Differential Geometry} for the K\"ahler manifold
background~\cite{NateDiffGeo}.

\subsection*{Elliptic Curves}

The advantage of scheme theory is in full effect when studying curves (Noetherian integral separated
$1$-dimensional schemes of finite type over $k$). The study of curves splits to
the study of curves of genus $g$ by the Riemann-Roch Theorem (see~\cite{NateComplexAnalysis}
or~\cite{NateAlgGeo}). Curves of genus $1$ are called \emph{elliptic curves} which by GAGA results
(\cite{hallGAGATheorems2022}) is equivalent to studying non-singular degree $3$ polynomial equations. When in
$\Char \ne 2,3$ we can algebraically say that the study of elliptic curves is the study of equation $y^2 = x^3
+ax + b$. Finding all rational solutions to such polynomials is still an open problem; the \emph{Birch
and Swinnerton-Dyer Conjecture} would, if proven, explain how many solutions there exists, and the person(s) who
succeeds to prove it would win \$1 million from the Clay Institute of Mathematics\footnote{It is one of the
Millennium problems set at the beginning of the 21st century}. Due to scheme
theory there are natural ways of understanding ``bad reductions'' (think of ramified primes), of defining
a covering space using galois theory (like in homotopy theory), and keeping track of when curves intersect
non-transversally (with multiplicity higher than $1$).

The reader interested can start on \emph{EYNTKA Elliptic Curves} (\cite{NateEllipticCurve}) or
look at \emph{Silvermen's Arithmetic of Elliptic Curves} (\cite{silvermanArithmeticEllipticCurves2009}).


\subsection*{Number Theory and Arithmetic Geometry}

Scheme theory over number fields or number rings is called \emph{Arithmetic Geometry}. A huge amount of
progress in modern theory was made in number theory by using tools in arithmetic geometry, for example the
proof of Fermat Last Theorem (that there exists no natural number $n\ge 3$ such that $x^n+y^n=y^n$ has
an integer solution) was proved through the lens of Arithmetic Geometry and analytic number theory.
Much of modern number theory builds directly on the algebraic foundations presented in this book.

A major area of research is the \emph{Langlands program} which aims to generalize Artin's reciprocity (see
\emph{Class Field Theory}, \cite{NateCFT}) which can be thought of the $1$-dimensional case of Langland's
reciprocity. The $2$-dimensional case strongly involved elliptic curves and the related modular forms and
automorphic forms. An introduction to the required background covering these
topics includes \cite{NateEllipticCurve}, \cite{NateAlgGeo}, and~\cite{silvermanAdvancedTopicsArithmetic1994}.

With the advent of complex analysis, many problems in number theory can be studied using certain meromorphic
functions. Of central interest is the \emph{Riemann-Zeta function}
\[
	\zeta(s) = \sum{n=1}^\infty \frac{1}{n^s} = \prod_{p\text{ prime}} \frac{1}{1-p^{-s}}
\]
which can be analytically extended to $\C\setminus\{1\}$ and captures an incredible amount of information
about the distribution of primes. The \emph{Riemann Hypothesis}, one of the Millennium problems, states that all the
[non-trivial]\footnote{There are roots that are given by the Gamma function which are considered to be related
to the prime at infinity, and is not important for the usual, finite, prime} roots of $\zeta(s)$ lie on
the line $L = \set{z \in \C}{\im(z) = 1/2}$. By some application of Fourier analysis, this would drastically
improve our understand of how many primes $p\le n$ there exists for any $n \in \N$, and would resolve many
open problems. Proving that all [nontrivial] zero's are on $L$ or that there is at least one [nontrivial] zero
that is not on $L$ would give the person(s) who proved it \$1 million by the Clay Institute.

\subsection*{Geometric Group Theory}

Infinite groups are much harder to grasp than finite groups. However there
are many infinite groups that naturally show up in geometric settings and have a metric that can be imposed
on them. Studying these groups falls under \emph{Geometric group theory} which views infinite
groups themselves as geometric objects. By equipping a group with a metric, one can study its large-scale
geometry to deduce its algebraic properties. This approach has led to the
resolution of major conjectures in topology (most famously, Thurston's Geometrization Conjecture implying the
\emph{Poincar\'e Conjecture}, another Millennium problem).


\subsection*{Derived Algebraic Geometry and Higher Algebra}

Schemes are incredibly power, but are not adept
with ``singular'' spaces or when taking ``non-regular'' quotients; categorically they do not behave well under
limit sand colimits. \emph{Derived geometry} is a
redevelopment of commutative algebra inspired abstract topology. It replaces rings with
derived objects that retain homological information, allowing for a more robust theory of
intersections and moduli spaces that works in greater generality. This field often draws on \emph{higher category theory} to create a unified language for
geometry.

\subsection*{Motivic and Chromatic Homotopy Theory}

Algebraic topology uses algebraic invariants, like homology groups, to classify and distinguish topological
spaces. The algebraic tools themselves have become the focus of intense study.
\emph{Chromatic homotopy theory} seeks to organize how shapes can be mapped
into one another (homotopy) by filtering it through algebraic data coming from the theory of formal groups.
\emph{Motivic homotopy theory} builds a direct bridge, or ``motive,'' between algebraic geometry and topology.
The ultimate goal is a structural understanding of shape itself, with the calculation of the
\emph{stable homotopy groups of spheres} remaining one of the field's major goals.

\subsection*{Anabelian Geometry}

Anabelian geometry is asking if we can recover an arithmetic variety $X$ given information about the (\'etale)
fundamental group of $X$. While  Motivic Homotopy Theory seeks to build a dictionary of algebraic invariants
to describe a space (using its entire tower of homotopy groups \( \pi_n \) of the sphere spectrum), anabelian
geometry exclusively on the first \'etale homotopy group, \( \pi_1^{\text{\'et}} \), and posits that for a
special class of ``hyperbolic'' varieties, this single algebraic object already contains all the necessary
information.


The central algebraic invariant in this field is the \emph{\'etale fundamental group},
denoted \( \pi_1^{\text{\'et}}(X) \). Its connection to the homotopy $\pi^\Top$ group comes from how the universal cover $E
\to X$ and the intermediate covers has a galois correspondence to the subgroups of $\pi^{\Top}$. An example of
a simple \'etale fundamental group is those for a field $k$, where $\pi^{\text{\'et}}(k) =\Gal_k(k_s)$ where
$k_s$ is the maximal separable extension of $k$.
More generally, for a variety \(X\) defined over a number field like \(
\mathbb{Q} \), its \'etale fundamental group is a profinite group that captures not only the
topological shape of \(X\) (as seen over the complex numbers) but also the entirety of its arithmetic
structure.

The ``anabelian'' hypothesis is the conjecture that this reconstruction is only possible when the group is sufficiently complex and far from being abelian.
\begin{itemize}
    \item \textbf{Abelian Case (Fails):} For an object like an elliptic curve, whose fundamental group is abelian, this principle fails. Many different elliptic curves can have the same fundamental group, so the group does not uniquely determine the geometry.
    \item \textbf{Anabelian Case (Conjectured to Succeed):} For a ``hyperbolic'' variety, such as a curve of
		genus two or more, the \'etale fundamental group is a non-abelian group. The anabelian hypothesis
		posits that there is enough information to uniquely determines the original variety.
\end{itemize}

This makes anabelian geometry a subfield of arithmetic geometry in its goals, but leaning more towards the
philosophy of geometric group theory -- it seeks to understand geometry by studying the algebraic properties
of an associated group. An unsolved problem (as of writing this book) is the \emph{Section Conjecture} which
attempts to use the structure of \( \pi_1^{\text{ét}}(X) \) to describe the set of rational points on \(X\),
offering a new (non-abelian) approach to solving Diophantine equations.



% \newpage
\printbibliography

\end{document}
